\chapter{Misspecification and the Edgeworth correction}
\label{ch:misspecification-edgeworth}

Two questions in the paper have the same mathematical shape: what does a maximum
likelihood estimator converge to when the model it fits is wrong, and what does the
Gaussian covariance twin lack when it stands in for the true score? Both are answered
by a projection. The first is the Kullback--Leibler projection of the true law onto the
model family, whose image is the \emph{pseudo-true parameter} of White; the paper turns
it into a closed-form asymptotic bias $b(\varepsilon)$ for each of its four
misspecification mechanisms (Section~5 of the paper, experiments E08 and E17). The second
is the expansion of the true noised law around its Gaussian twin in Hermite polynomials,
whose first non-trivial term is the fourth-cumulant correction of Edgeworth; the paper
adds that term back and measures how much of the closure gap it closes (Section~3 of the
paper, experiment E16). This chapter builds both, derives the formulas the paper quotes,
and records how the experiments used them.

\section{The Kullback--Leibler projection and the pseudo-true parameter}
\label{sec:kl-projection}

\subsection{What a misspecified likelihood converges to}

Let the data $\vx^1,\dots,\vx^{\nseq}$ be i.i.d.\ from a law $P^\star$ and let
$\{Q_\params\}$ be a family of densities that may not contain $P^\star$. The
\term{misspecified maximum-likelihood estimator} is
$\hat\params=\argmax_\params\sum_\mu\log Q_\params(\vx^\mu)$. Dividing by $\nseq$, the
law of large numbers gives
\begin{equation}
  \frac1\nseq\sum_\mu\log Q_\params(\vx^\mu)\ \to\ \E_{P^\star}\log Q_\params(\vx)
  =\E_{P^\star}\log P^\star(\vx)-\KL\bigl(P^\star\,\|\,Q_\params\bigr),
  \label{eq:pop-loglik}
\end{equation}
since $\KL(P^\star\|Q)=\E_{P^\star}[\log P^\star-\log Q]$ (\Cref{ch:conventions}). The first
term does not depend on $\params$.

\begin{definition}[Pseudo-true parameter; Kullback--Leibler projection]
\label{def:pseudo-true}
The \term{pseudo-true parameter} is
\begin{equation}
  \params^\ast=\argmin_\params\KL\bigl(P^\star\,\|\,Q_\params\bigr),
  \label{eq:pseudo-true}
\end{equation}
and $Q_{\params^\ast}$ is the \term{Kullback--Leibler projection} of $P^\star$ onto the
family. When $P^\star=Q_{\params_0}$ belongs to the family, $\params^\ast=\params_0$.
\end{definition}

\begin{proposition}[White]
\label{prop:white}
If the minimiser \cref{eq:pseudo-true} is unique and the usual regularity holds, then
$\hat\params\to\params^\ast$ in probability, and
$\sqrt\nseq(\hat\params-\params^\ast)$ is asymptotically normal with the sandwich
covariance $\mat{J}^{-1}\mat{V}\mat{J}^{-1}$, $\mat{J}=-\E_{P^\star}\partial^2_\params\log Q_{\params^\ast}$,
$\mat{V}=\Var_{P^\star}(\partial_\params\log Q_{\params^\ast})$.
\end{proposition}

The proof is the M-estimator argument of \Cref{ch:quasi-likelihood-moments}: the
objective converges uniformly to \cref{eq:pop-loglik}, whose unique maximiser is
$\params^\ast$. The content of the proposition is the reading: a wrong model does not
make the estimator wander; it makes it converge, to a definite point, and more data
only sharpens the convergence. The error that remains at $\params^\ast$ is a
\term{bias} in the strict sense that no sample size removes it. This is the sentence
the paper's misspecification section ends on.

\subsection{Why EM--BP's limit is a Gaussian projection here}

The paper's misspecification experiment fits the mixture-innovation AR(1) kernel of
\Cref{ch:em-through-noise} to data from a Gaussian truth that is not an AR(1) chain. The
fitted innovation turns out Gaussian to three decimals (the paper reports the fitted
excess kurtosis; E08 records it per cell), which is what one expects: with a Gaussian
truth there is no fourth-order structure for the mixture to fit, and the extra
components only add variance. The paper therefore takes the $\nseq\to\infty$ limit of
EM--BP to be the projection onto the \emph{Gaussian} AR(1) family, a two-parameter family
$(\corr,\ivar)$, and computes it in closed form. Strictly the projection onto the
mixture family could be non-Gaussian even for a Gaussian truth, since the truth is
outside both families; the paper's justification is empirical (the fitted kurtosis)
and the closed form is checked against the measured $\nseq=2048$ errors, which is the
right order of argument.

\section{The paper's closed-form bias}
\label{sec:bias}

\subsection{The projection under the one-view protocol}

Let the truth be Gaussian with clean covariance $\Sig(\varepsilon)$, so its noised law at
level $t$ is $\normal{0}{\Sig_t(\varepsilon)}$ with $\Sig_t(\varepsilon)=e^{-2t}\Sig(\varepsilon)+\Dt\Id$.
The model is the Gaussian chain with $\hat\Sig_t(\corr,\ivar)=e^{-2t}\Sig_0(\corr,\ivar)+\Dt\Id$,
where $\Sig_0(\corr,\ivar)$ is the general chain covariance of
\Cref{ch:quasi-likelihood-moments}. Under the one-view protocol the $\nseq$ sequences are
split evenly over the twelve levels, so the population log-likelihood is an average
over levels with equal weights, and the pseudo-true parameter is
\begin{equation}
  \params^\ast(\varepsilon)=\argmin_{(\corr,\ivar)}\sum_{t\in\mathcal T}
  \KL\Bigl(\normal{0}{\Sig_t(\varepsilon)}\,\Big\|\,\normal{0}{\hat\Sig_t(\corr,\ivar)}\Bigr),
  \label{eq:theta-star}
\end{equation}
which is the formula of the paper's Appendix A5. The Gaussian divergence is
\cref{eq:kl-gauss} of \Cref{ch:quasi-likelihood-moments}, so up to terms without
$\params$ the objective is
\begin{equation}
  F(\corr,\ivar)=\frac12\sum_t\Bigl[\log\det\hat\Sig_t+\tr\bigl(\hat\Sig_t^{-1}\Sig_t(\varepsilon)\bigr)\Bigr],
  \label{eq:F}
\end{equation}
exactly what \code{kl\_project} in \code{experiments/e17\_misspec\_bias.py} minimises by
Nelder--Mead over $(\tanh^{-1}\corr,\log\ivar)$, the reparameterisation keeping
$|\corr|<1$ and $\ivar>0$. Setting the gradient of $F$ to zero gives the two
\term{projection equations}
\begin{equation}
  \sum_t\tr\Bigl(\hat\Sig_t^{-1}\,\partial_\theta\hat\Sig_t\,\hat\Sig_t^{-1}\bigl(\Sig_t(\varepsilon)-\hat\Sig_t\bigr)\Bigr)=0,
  \qquad\theta\in\{\corr,\ivar\},
  \label{eq:proj-eqs}
\end{equation}
the population version of the quasi-score equations of
\Cref{ch:quasi-likelihood-moments}: the projected covariance agrees with the truth in the
two directions $\hat\Sig_t^{-1}\partial_\theta\hat\Sig_t\hat\Sig_t^{-1}$, summed over the levels.

\subsection{The relative score error of the projected model}

At the projection the model's score is $-\hat\Sig_t^{-1}\vx$ and the truth's is
$-\Sig_t^{-1}\vx$ (write $\Sig_t$ for $\Sig_t(\varepsilon)$). The paper's estimand pools
squared errors and reference energies over the schedule
(\Cref{ch:experimental-statistics}); its population version is
\begin{equation}
  b(\varepsilon)^2=\frac{\sum_t\E\lVert(\hat\Sig_t^{-1}-\Sig_t^{-1})\vx\rVert^2}{\sum_t\E\lVert\Sig_t^{-1}\vx\rVert^2}
  \qquad(\vx\sim\normal{0}{\Sig_t}\text{ at level }t),
\end{equation}
with equal numbers of sequences at every level, as in the test bundles. Both
expectations are of quadratic forms: for $\vx\sim\normal0{\Sig}$ and any matrix $\mat{D}$,
$\E\lVert\mat{D}\vx\rVert^2=\E[\vx^\top\mat{D}^\top\mat{D}\vx]=\tr(\mat{D}^\top\mat{D}\Sig)=\tr(\mat{D}\Sig\mat{D}^\top)$.
With $\mat{D}=\hat\Sig_t^{-1}-\Sig_t^{-1}$ (symmetric) in the numerator and
$\mat{D}=\Sig_t^{-1}$ in the denominator,
\begin{equation}
  b(\varepsilon)^2=\frac{\sum_t\tr\bigl[(\hat\Sig_t^{-1}-\Sig_t^{-1})\,\Sig_t\,(\hat\Sig_t^{-1}-\Sig_t^{-1})\bigr]}{\sum_t\tr\,\Sig_t^{-1}},
  \qquad\hat\Sig_t=\hat\Sig_t\bigl(\params^\ast(\varepsilon)\bigr),
  \label{eq:bias}
\end{equation}
which is the trace formula of the paper's Appendix A5 and the function
\code{pooled\_err} of E17. No Monte Carlo enters: $b(\varepsilon)$ is a deterministic
function of the mechanism, its strength, the chain length and the schedule.

E17 computes a second projection beside the Kullback--Leibler one: the \term{lag-one
plug-in}, the AR(1) covariance $\corr_1^{|i-j|}$ whose correlation $\corr_1$ is the mean
lag-one correlation of the truth. It is the population limit of the moment estimator of
\Cref{ch:quasi-likelihood-moments} and of E08's \code{lmmse\_ar1} arm, and the paper
quotes both, noting that the measured EM--BP fit sits between the two because the
mixture kernel is not exactly Gaussian and a finite-sample fit undershoots the
Kullback--Leibler optimum.

\subsection{The finite-sample prediction \texorpdfstring{$\sqrt{b^2+v/\nseq}$}{sqrt(b2 + v/M)}}

At finite $\nseq$ the fitted parameter is $\hat\params=\params^\ast+\delta$ with
$\delta=O(\nseq^{-1/2})$ by \Cref{prop:white}. Expanding the model score to first order
in $\delta$, $\hat s=s_{\params^\ast}+\delta\cdot\partial_\params s_{\params^\ast}$, and
the squared error against the truth is
\begin{equation}
  \lVert\hat s-s\rVert^2=\lVert s_{\params^\ast}-s\rVert^2
  +2\,\delta\cdot\bigl\langle\partial_\params s_{\params^\ast},\,s_{\params^\ast}-s\bigr\rangle
  +\lVert\delta\cdot\partial_\params s_{\params^\ast}\rVert^2 .
\end{equation}
Take expectations over the test data and over the training sample that produced
$\delta$. The first term is $b^2$ times the reference energy. The cross term vanishes to
leading order because $\E\delta=o(\nseq^{-1/2})$ and $\delta$ is independent of the test
sequences. The last term is a quadratic form in $\delta$ whose expectation is
$\tr(\Var(\delta)\,\mat{G})$ for a fixed matrix $\mat{G}$, hence $v/\nseq$ for a constant
$v$ that depends on the sandwich covariance and the mechanism. Therefore
\begin{equation}
  \text{risk}(\nseq,\varepsilon)^2\approx b(\varepsilon)^2+\frac{v(\varepsilon)}{\nseq}.
  \label{eq:finite-sample}
\end{equation}
The paper reads $v$ off the well-specified cells: at $\varepsilon=0$ the bias vanishes,
so $v=\nseq\cdot\text{risk}(\nseq,0)^2$, and it takes $v$ to be the same at every
strength (E17's setting \code{v\_const}), a simplification that the agreement at both
sizes in every mechanism justifies after the fact. The two sizes of E08, $512$ and $2048$,
separate the two terms: where the risk falls by a factor of two between them the
variance term dominates, and where it does not move the bias does. The paper's Figure~5
(bottom row) draws $b(\varepsilon)$ and \cref{eq:finite-sample} through the measured
points (\Cref{ch:reading-the-figures}).

\section{The order of the bias in the strength}
\label{sec:bias-order}

\subsection{First-order perturbation of a projection}

Write $\Sig(\varepsilon)=\Sig_0+\varepsilon\mat{D}+O(\varepsilon^2)$ with $\Sig_0=\Sig_0(\params_0)$
in the family, and consider one level for simplicity (the sum over levels adds the
same terms). Linearising the projection equations \cref{eq:proj-eqs} around
$\params_0$ gives $\params^\ast(\varepsilon)=\params_0+\varepsilon\,\delta\params+O(\varepsilon^2)$ with
\begin{equation}
  \delta\params=\mat{I}(\params_0)^{-1}\,\bigl\langle\Sig_0^{-1}\partial_\theta\Sig_0\Sig_0^{-1},\mat{D}\bigr\rangle_\theta,
  \qquad
  \mat{I}_{\theta\theta'}=\bigl\langle\Sig_0^{-1}\partial_\theta\Sig_0\Sig_0^{-1},\partial_{\theta'}\Sig_0\bigr\rangle,
\end{equation}
where $\langle\mat{A},\mat{B}\rangle=\tr(\mat{A}\mat{B})$: the projection moves the
parameter by the \emph{Fisher-metric projection} of the perturbation $\mat{D}$ onto the
tangent space of the family, spanned by $\partial_\corr\Sig_0$ and $\partial_\ivar\Sig_0$. The
projected covariance is then
$\hat\Sig(\varepsilon)=\Sig_0+\varepsilon\,\delta\params\cdot\partial_\params\Sig_0+O(\varepsilon^2)$, and
the mismatch $\hat\Sig-\Sig=\varepsilon(\delta\params\cdot\partial_\params\Sig_0-\mat{D})+O(\varepsilon^2)$
is $\varepsilon$ times the component of $\mat{D}$ \emph{orthogonal} to the tangent space in
the Fisher metric. Inserting into \cref{eq:bias}, with
$\hat\Sig^{-1}-\Sig^{-1}=-\Sig^{-1}(\hat\Sig-\Sig)\Sig^{-1}+O(\varepsilon^2)$,
\begin{equation}
  b(\varepsilon)=\varepsilon\,\lVert\mat{D}^\perp\rVert_{\Sig_0^{-1}}\big/\sqrt{\tr\Sig_0^{-1}}+O(\varepsilon^2),
  \qquad \lVert\mat{A}\rVert^2_{\Sig^{-1}}=\tr(\Sig^{-1}\mat{A}\Sig^{-1}\mat{A}\Sig^{-1}\Sig)=\tr\bigl((\Sig^{-1}\mat{A})^2\bigr).
  \label{eq:bias-first-order}
\end{equation}

\begin{proposition}[Order of the bias]
\label{prop:bias-order}
$b(\varepsilon)=\Theta(\varepsilon)$ if the first-order perturbation $\mat{D}$ of the
covariance has a nonzero component outside the tangent space of the family at
$\params_0$, and $b(\varepsilon)=\Theta(\varepsilon^2)$ if $\mat{D}=0$, that is, if the
covariance depends on the strength only through $\varepsilon^2$.
\end{proposition}

The second case is the \term{rank-one latent} of E08: $a=(y+\beta g)/\sqrt{1+\beta^2}$
with $y$ an AR(1) chain and $g$ a single standard normal shared by all sites, so
$\Cov(a_i,a_j)=(\corr^{|i-j|}+\beta^2)/(1+\beta^2)$ for $i\ne j$ (\code{GaussianAR1PlusGlobal}
in \code{src/markov\_score/priors.py}). The covariance is even in $\beta$, its first-order
term in $\beta$ vanishes, and the bias is $\Theta(\beta^2)$. The audit adds a second
reason the latent is weak: the perturbation $\beta^2\mathbf 1\mathbf 1^\top$ has rank one,
so it moves one eigendirection of $\Sig_0$ among $\nsites$, and its share of the score
energy is $O(1/\nsites)$ as well as $O(\beta^2)$. The other three mechanisms perturb
of order $\nsites$ directions at first order and have $b=\Theta(\varepsilon)$, which is
what E17's small-$\varepsilon$ local exponents report (the paper quotes them as macros).

\begin{example}[A projection and its bias, by hand]
\label{ex:projection}
Take $\nsites=2$ and a single level with $t\to0$, so that $\Sig_t\to\Sig$ and the
model covariance is $\hat\Sig(\corr,\ivar)=\begin{psmallmatrix}1&\corr\\ \corr&\corr^2+\ivar\end{psmallmatrix}$,
with $\det\hat\Sig=\ivar$ and
$\hat\Sig^{-1}=\ivar^{-1}\begin{psmallmatrix}\corr^2+\ivar&-\corr\\-\corr&1\end{psmallmatrix}$.
Let the truth be $\Sig(\varepsilon)=\begin{psmallmatrix}1+\varepsilon&c\\ c&1-\varepsilon\end{psmallmatrix}$,
a chain with lag-one correlation $c$ whose two variances have been tilted by
$\pm\varepsilon$; this perturbation lies outside the family, whose first variance is
fixed at $1$. The objective \cref{eq:F} is
$F=\ivar^{-1}[(\corr^2+\ivar)(1+\varepsilon)-2\corr c+1-\varepsilon]+\log\ivar$.
Setting $\partial_\corr F=0$ gives $\corr^\ast=c/(1+\varepsilon)$; setting
$\partial_\ivar F=0$ gives $\ivar^\ast=\corr^{\ast2}(1+\varepsilon)-2\corr^\ast c+1-\varepsilon=1-\varepsilon-c^2/(1+\varepsilon)$.
The projected covariance is
\begin{equation}
  \hat\Sig^\ast=\begin{pmatrix}1&c/(1+\varepsilon)\\ c/(1+\varepsilon)&1-\varepsilon-c^2\varepsilon/(1+\varepsilon)^2\end{pmatrix}
  =\Sig(\varepsilon)-\varepsilon\,\vect{v}\vect{v}^\top+O(\varepsilon^2),\qquad\vect{v}=\begin{pmatrix}1\\ c\end{pmatrix}.
\end{equation}
So the residual mismatch is rank one, along $\vect{v}$: the family absorbs the
perturbation of the second variance and of the lag-one entry, and what it cannot absorb
is the first variance, which it pins at one. Then
$\hat\Sig^{\ast-1}-\Sig^{-1}=\varepsilon\,\Sig^{-1}\vect{v}\vect{v}^\top\Sig^{-1}+O(\varepsilon^2)$ and
the numerator of \cref{eq:bias} is $\varepsilon^2(\vect{v}^\top\Sig^{-1}\vect{v})^2$. At
$\varepsilon=0$, $\Sig_0^{-1}=(1-c^2)^{-1}\begin{psmallmatrix}1&-c\\-c&1\end{psmallmatrix}$ and
$\vect{v}^\top\Sig_0^{-1}\vect{v}=(1-2c^2+c^2)/(1-c^2)=1$, while the denominator is
$\tr\Sig_0^{-1}=2/(1-c^2)$. Hence
\begin{equation}
  b(\varepsilon)=\varepsilon\sqrt{\frac{1-c^2}{2}}+O(\varepsilon^2),
\end{equation}
linear in the strength as \Cref{prop:bias-order} says, with a coefficient that is
smaller for a more correlated chain: at $c=0.85$, $\sqrt{0.2775/2}\approx0.37$, so a
ten per cent tilt of the variances costs a relative score error of about $0.04$.
Had the truth been $\begin{psmallmatrix}1&c+\varepsilon\\ c+\varepsilon&1\end{psmallmatrix}$
instead, the perturbation would lie in the family (it is the chain at correlation
$c+\varepsilon$ with $\ivar=1-(c+\varepsilon)^2$) and $b\equiv0$.
\end{example}

\section{The misspecification mechanisms of E08}
\label{sec:e08}

E08 (\code{experiments/e08\_misspecification.py}) keeps every truth Gaussian, so that its
score is $-\Sig_t^{-1}\vx$ by linear algebra and no grid enters the reference. The four
mechanisms, each with a strength whose zero is the well-specified chain at
$\corr=0.85$, $\nsites=32$:
\begin{description}[leftmargin=2em,itemsep=0.25em]
\item[\code{global}, rank-one latent $\beta$.] $\Cov(a_i,a_j)=(\corr^{|i-j|}+\beta^2)/(1+\beta^2)$
  off the diagonal, unit variances; Gaussian and variance-preserving but not Markov.
  Strengths $\beta\in\{0,0.1,0.25,0.5,1\}$.
\item[\code{longrange}, long-range precision coupling $\gamma$.] The precision matrix is
  the tridiagonal AR(1) precision plus $\gamma$ times a matrix with entries
  $-e^{-|i-j|/8}$ for $|i-j|\ge2$ and a diagonal compensation that keeps it diagonally
  dominant; the covariance is its inverse rescaled to unit diagonal
  (\code{GaussianLongRange}). Strengths $\gamma\in\{0,0.05,0.1,0.2,0.4\}$.
\item[\code{ar2}, a second autoregressive coefficient $\corr_2$.] A stationary AR(2)
  $a_i=\corr_1a_{i-1}+\corr_2a_{i-2}+\varepsilon_i$ with $\corr_1=\corr(1-\corr_2)$, which keeps
  the lag-one correlation at $\corr$ (from the first Yule--Walker equation
  $\gamma_1=\corr_1/(1-\corr_2)$); the covariance follows from the Yule--Walker
  recursion $\gamma_k=\corr_1\gamma_{k-1}+\corr_2\gamma_{k-2}$ (\code{\_ar2\_cov}).
  Strengths $\corr_2\in\{0,0.1,0.2,0.3,0.4\}$.
\item[\code{sitevar}, site-dependent correlation $\delta$.] $a_i=\corr_ia_{i-1}+\varepsilon_i$
  with $\corr_i=\corr+\delta\cos(2\pi i/\nsites)$ and $\Var\varepsilon_i=1-\corr_i^2$, so every
  site keeps unit variance and $\Cov(a_i,a_j)=\prod_{k=i+1}^{j}\corr_k$
  (\code{\_sitevar\_cov}). The construction needs $|\corr_i|<1$ at every site, which at
  $\corr=0.85$ bounds the amplitude by $0.15$; the registered strength $0.2$ was invalid
  by construction and was replaced by $0.14$ in a recorded amendment before any number
  was read (\Cref{ch:protocol}). Strengths $\delta\in\{0,0.05,0.1,0.14\}$.
\end{description}
Each truth is fitted by four arms on identical bundles: \code{em\_bp} (the
eight-component mixture kernel, one noised view per sequence, checkpoint on the
validation bulk risk); \code{window} (a radius-$2$, width-$64$ window head trained under
interface~A for $16{,}000$ steps with the ladder $\{500,\dots,16000\}$, checkpoint on the
validation bulk risk; \Cref{ch:neural-baselines}); \code{lmmse\_ar1} (the LMMSE of the
AR(1) covariance with the truth's mean lag-one correlation: the best a Gaussian AR(1)
with that correlation can do); and \code{isotropic}. Two sizes, $\nseq\in\{512,2048\}$, so
that the estimation term and the bias term of \cref{eq:finite-sample} can be told apart.
The predeclared analysis (\code{tools/make\_misspec\_numbers.py}) is the mean and standard
error of each arm per cell, the paired per-seed ratio window/EM--BP, and the
\term{break-even} strength: the paper reports the two grid points that bracket the
sign change of the mean log-ratio and never a fitted crossing, because with the
available seeds the crossing is resolved only to the grid. The registry entry declares
eight seeds; the script's default is the protocol's seed list, and the paper's macro for
the seed count records what was actually run.

\section{The Edgeworth correction around the covariance twin}
\label{sec:edgeworth}

\subsection{Expanding a law around a Gaussian with the same covariance}

Let $\vx\sim\Pt$ have mean zero, covariance $\Sig_t$ and fourth-cumulant tensor
$\mat{K}=\kfour[\vx]=e^{-4t}\kfour[\va]$ (\Cref{ch:attenuation-theorems}); its third
cumulants vanish by the symmetry of the chain. The characteristic function of $\Pt$
and of its twin $\normal0{\Sig_t}$ have log-characteristic functions that agree to second
order and differ at fourth:
\begin{equation}
  \log\E e^{i\vect{s}\cdot\vx}=-\tfrac12\vect{s}^\top\Sig_t\vect{s}+\tfrac{1}{24}\mat{K}[\vect{s},\vect{s},\vect{s},\vect{s}]+\dots
  \qquad(\text{cumulant expansion, \Cref{ch:moments-cumulants}}).
\end{equation}
Exponentiating and keeping the first correction,
$\E e^{i\vect{s}\cdot\vx}\approx e^{-\frac12\vect{s}^\top\Sig_t\vect{s}}\bigl(1+\tfrac1{24}\mat{K}[\vect{s},\vect{s},\vect{s},\vect{s}]\bigr)$.
Under the inverse Fourier transform a factor $(i s_j)$ becomes $-\partial_j$, so
$s_is_js_ks_l$ becomes $\partial_i\partial_j\partial_k\partial_l$ acting on the Gaussian density
$\phi_{\Sig_t}$:
\begin{equation}
  \Pt(\vx)\approx\phi_{\Sig_t}(\vx)\Bigl[1+\frac1{24}\sum_{ijkl}K_{ijkl}\,H^{\Sig_t}_{ijkl}(\vx)\Bigr],
  \qquad
  H^{\Sig}_{ijkl}(\vx)=\frac{\partial_i\partial_j\partial_k\partial_l\phi_\Sig(\vx)}{\phi_\Sig(\vx)} .
  \label{eq:edgeworth}
\end{equation}
This is the \term{first Edgeworth correction} (also called the Gram--Charlier
correction at this order; the two series differ in how they group higher terms, not in
this term). The tensors $H^\Sig_{ijkl}$ are the \term{multivariate Hermite polynomials}
of the Gaussian $\normal0\Sig$. With $\vect{y}=\Sig^{-1}\vx$ one has
$\partial_i\phi_\Sig=-y_i\phi_\Sig$ and $\partial_iy_j=(\Sig^{-1})_{ij}$, and four
differentiations give
\begin{equation}
  H^\Sig_{ijkl}=y_iy_jy_ky_l-\bigl[y_iy_j(\Sig^{-1})_{kl}+\text{five more pairings}\bigr]
  +\bigl[(\Sig^{-1})_{ij}(\Sig^{-1})_{kl}+\text{two more}\bigr],
\end{equation}
the exact analogue of $\Hen{4}(x)=x^4-6x^2+3$ (\Cref{ch:gaussian-hermite}). Contracting
with the symmetric tensor $\mat{K}$ collapses the six pairings into one term and the
three into one:
\begin{equation}
  \epsilon(\vx):=\frac1{24}\sum_{ijkl}K_{ijkl}H^{\Sig_t}_{ijkl}(\vx)
  =\frac1{24}\Bigl[\mat{K}[\vect{y},\vect{y},\vect{y},\vect{y}]-6\,\vect{y}^\top\mat{A}\vect{y}+3\,\langle\mat{A},\Sig_t^{-1}\rangle\Bigr],
  \label{eq:eps}
\end{equation}
with $A_{ml}=\sum_{jk}K_{mjkl}(\Sig_t^{-1})_{jk}$ a symmetric matrix (the partial
contraction of the tensor with the precision), which is the formula of the paper's
Appendix A5 and of E16.

\subsection{The corrected score}

The score of the corrected density $\phi_{\Sig_t}(1+\epsilon)$ is
\begin{equation}
  \score^{\mathrm{E4}}(\vx)=\nabla\log\phi_{\Sig_t}+\nabla\log(1+\epsilon)=\Sgauss(\vx)+\frac{\nabla\epsilon(\vx)}{1+\epsilon(\vx)},
  \label{eq:score-e4}
\end{equation}
with $\Sgauss=-\Sig_t^{-1}\vx$ the twin's score, and, differentiating \cref{eq:eps} with
$\partial_my_j=(\Sig_t^{-1})_{mj}$,
\begin{equation}
  \nabla\epsilon=\frac16\,\Sig_t^{-1}\mat{K}[\,\cdot\,,\vect{y},\vect{y},\vect{y}]-\frac12\,\Sig_t^{-1}\mat{A}\vect{y}
  \label{eq:grad-eps}
\end{equation}
(the quartic gives four equal terms, the quadratic two, and the constant none). E16
evaluates both the \term{log form} \cref{eq:score-e4} and the \term{linear form}
$\Sgauss+\nabla\epsilon$, which drops the denominator; the two agree where $\epsilon$ is
small and differ where it is not.

\subsection{When the expansion diverges}

\Cref{eq:edgeworth} is an asymptotic expansion in the size of the fourth cumulant
relative to the covariance. Its correction $\epsilon(\vx)$ is a quartic in
$\vect{y}$ with coefficient of order $e^{-4t}\Frob{\kfour[\va]}$, so three things go
wrong at low noise. The correction is not small on the bulk of the distribution, so the
truncated series is not close to $\Pt$; $1+\epsilon$ can be negative, so the corrected
``density'' is not one and the log form's denominator can change sign (E16 clips it at
$0.1$); and the quartic always wins in the tail, so far out the correction is large at
every $t$. The expansion is useful exactly where the paper's Theorem~3 is an asymptotic
statement, above the crossover, and the paper says so: below a level among the lowest
rungs of the schedule (the paper quotes it as a macro) the correction hurts, and pooled over the twelve-level schedule the corrected score is
worse than the twin because the lowest levels carry most of the reference energy.

\begin{example}[The scalar correction, by hand]
\label{ex:edgeworth-scalar}
For $\nsites=1$ and unit variance, $\Sig_t=1$, $\vect{y}=x$, $\mat{K}=\mat{A}=K$, and
\cref{eq:eps} reads $\epsilon(x)=\tfrac{K}{24}(x^4-6x^2+3)=\tfrac{K}{24}\Hen{4}(x)$, with
$\nabla\epsilon=\tfrac K6(x^3-3x)=\tfrac K6\Hen{3}(x)$: the scalar leading term of the
paper's Theorem~3, whose $\gamma$-second moment $\tfrac{K^2}{36}\E\Hen{3}^2=\tfrac{K^2}{36}\cdot3!=K^2/6$
is the closure constant. For a unit-variance Laplace law $\kfour=3$, so $K=3e^{-4t}$.
At $t=1$, $K=3e^{-4}\approx0.0549$ and $\epsilon(x)\approx0.00229\,\Hen{4}(x)$: at
$x=0$ it is $0.0069$, at $x=2$ ($\Hen{4}=-5$) it is $-0.011$, at $x=3$ ($\Hen{4}=30$) it
is $0.069$, all small, and the correction is a genuine improvement of the twin over
the bulk; at $x=5$ ($\Hen{4}=478$) it is $1.09$, larger than one, but $x=5$ is five
standard deviations out. At $t=0.1$, $K=3e^{-0.4}\approx2.01$ and
$\epsilon(x)\approx0.0838\,\Hen{4}(x)$: $0.25$ at $x=0$, $-0.42$ at $x=2$, $2.5$ at
$x=3$. The correction is of order one on the bulk and the series is not an expansion
at all there. The dividing line, where $\tfrac{K}{24}\max_{|x|\le2}|\Hen{4}|=\tfrac{5K}{24}$
reaches, say, one tenth, is $K\approx0.5$, that is $e^{-4t}\approx1/6$, $t\approx0.45$; on
the chain the relevant size is $e^{-4t}\Frob{\kfour[\va]}$ against $\lmax(\Sig_0)^2$,
which puts the line near the crossover $\tstar$, consistent with what E16 measures.
\end{example}

\subsection{How E16 used it}

E16 (\code{experiments/e16\_edgeworth\_oracle.py}) takes the closure protocol of E01:
$512$ chains from the stream \code{rng\_for("e01-closure", family, 0)} for the Laplace,
bimodal, uniform and Student-$t_5$ chains at $\corr=0.85$, one common Gaussian draw
$\vz$ reused across the twelve levels and the two extra levels $3.5$ and $4$, and the
tree-exact reference on the $(801,8)$ grid. Per family and level it computes the twin's
score, the linear and the log corrected scores with the \emph{true} $\Sig_0$ and the
\emph{true} tensor $\kfour[\va]$ (\code{theory.chain\_kappa4}, an oracle), and reports
the relative score error of each against the reference and the fraction of the squared
closure gap closed,
\begin{equation}
  \text{closed}(t)=1-\frac{\sum_j\lVert\score^{\mathrm{E4}}(\vx_j)-\score(\vx_j)\rVert^2}{\sum_j\lVert\Sgauss(\vx_j)-\score(\vx_j)\rVert^2},
\end{equation}
which is one when the correction is the whole gap, zero when it does nothing, and
negative when it hurts. The predeclared reading is the level above which the log form
closes more than ninety per cent, the level below which it is negative, and the pooled
error of the three estimators over the schedule. The paper quotes the closed fractions at
several levels for the Laplace chain and the ninety-per-cent levels for the bimodal and
uniform chains (Section~3 and Appendix A5, macros \code{audEdge\ldots}). Two cautions
the registry and the audit attach: on Student-$t_5$ the reference on the $(801,8)$ grid
is dominated by the truncated prior above $t\approx3$ (\Cref{ch:grid-quadrature}), so the
Student-$t$ fractions at those levels measure the truncation and are not used; and the
oracle uses the true fourth cumulant, so it bounds what a moment method built on
fourth-order statistics could do rather than being one (a version with the moment
estimate of \Cref{ch:quasi-likelihood-moments} in place of the true tensor would be the
competitor, and the audit recommends it).

\where{\begingroup\sloppy The paper, Section~3 (``the gap is the fourth-order term''),
Section~5 (misspecification; Figure~5) and Appendix A5 (E16, E17). Code:
\code{experiments/e17\_misspec\_bias.py} (\code{kl\_project}, \code{pooled\_err}),
\code{experiments/e08\_misspecification.py} (\code{\_truth}, \code{\_ar2\_cov},
\code{\_sitevar\_cov}), \code{src/markov\_score/priors.py} (\code{GaussianAR1PlusGlobal},
\code{GaussianLongRange}), \code{experiments/e16\_edgeworth\_oracle.py},
\code{src/markov\_score/theory.py} (\code{chain\_kappa4}, \code{score\_gap\_constant}),
\code{tools/make\_misspec\_numbers.py}, \code{tools/make\_audit\_numbers.py}.\par\endgroup}

\paragraph{References.}
\begingroup\sloppy
H.~White, 1982, Maximum likelihood estimation of misspecified models, \emph{Econometrica} 50, 1--25.
S.~Kullback and R.~A.~Leibler, 1951, On information and sufficiency, \emph{Annals of Mathematical Statistics} 22, 79--86.
H.~Cram\'er, 1946, \emph{Mathematical Methods of Statistics}, Princeton University Press (chapter 17: the Edgeworth and Gram--Charlier series).
P.~McCullagh, 1987, \emph{Tensor Methods in Statistics}, Chapman and Hall (multivariate Hermite polynomials and cumulant tensors).
O.~E.~Barndorff-Nielsen and D.~R.~Cox, 1989, \emph{Asymptotic Techniques for Use in Statistics}, Chapman and Hall.
A.~Katsevich and A.~S.~Bandeira, 2023, Likelihood maximization and moment matching in low SNR Gaussian mixture models, \emph{Communications on Pure and Applied Mathematics} 76, 788--842.\par\endgroup
